\begin{table}[htbp]
\centering
\small
\caption{Problema E, Passos 3 e 6 (google/gemini-2.5-flash-lite): LLM com 40 exemplos versus classificadores clássicos.}
\label{tab:bloco2_E_passo6_baselines_flashlite}
\begin{tabular}{lrlrrr}
\toprule
 & \multicolumn{5}{c}{Acur.\ vs espec.\ (\%)} \\ \cmidrule(lr){2-6}
Especialista & LLM (mista) & LLM (melhor estratégia) & Regressão Logística & SVM (RBF) & k-NN (k=5) \\
\midrule
E1 & 97,6 & 97,6 (Mista) & 100,0 & 95,2 & 91,2 \\
E2 & 98,8 & 100,0 (Difícil) & 100,0 & 95,8 & 93,3 \\
E3 & 93,9 & 95,4 (Aleatória) & 100,0 & 99,4 & 98,2 \\
\bottomrule
\multicolumn{6}{p{0.95\linewidth}}{\footnotesize Problema E, Passos 3 (LLM) e 6 (baselines). Acurácia (LLM vs especialista), em \%, com k = 40 exemplos (maior k da grade). Os classificadores clássicos (k-NN, Regressão Logística, SVM-RBF) são treinados nos MESMOS k exemplos da estratégia mista e avaliados no MESMO conjunto de teste; não dependem do LLM. Coluna LLM (mista): mesma estratégia dos baselines.} \\
\end{tabular}
\end{table}
